% BHU PYQ -- Professional Exam Template
\documentclass[11pt,a4paper]{article}

\usepackage[a4paper,top=2.5cm,bottom=2.5cm,left=2.8cm,right=2.8cm,headheight=14pt]{geometry}
\usepackage{amsmath,amssymb}
\usepackage[T1]{fontenc}
\usepackage[utf8]{inputenc}
\usepackage{lmodern}
\usepackage{microtype}
\usepackage{fancyhdr}
\usepackage{enumitem}
\usepackage{listings}
\usepackage{hyperref}
\usepackage{lastpage}
\usepackage{parskip}
\usepackage{booktabs}

\hypersetup{colorlinks=true,urlcolor=black,linkcolor=black,
  pdftitle={CS-106: Operating System Concepts -- BHU PYQ}}

\pagestyle{fancy}
\fancyhf{}
\renewcommand{\headrulewidth}{0.4pt}
\renewcommand{\footrulewidth}{0.4pt}
\fancyhead[L]{\small\textit{Banaras Hindu University}}
\fancyhead[R]{\small\textit{CS-106: Operating System Concepts}}
\fancyfoot[C]{\small Page \thepage\ of \pageref{LastPage}}

\lstset{
  basicstyle=\small\ttfamily,
  language=C,
  frame=single,
  breaklines=true,
  columns=flexible,
  keepspaces=true,
  numbers=left,
  numberstyle=\tiny,
  showstringspaces=false
}

\newlist{parts}{enumerate}{2}
\setlist[parts,1]{label=(\alph*),leftmargin=*,itemsep=4pt,topsep=3pt}
\setlist[parts,2]{label=(\roman*),leftmargin=*,itemsep=2pt,topsep=2pt}
\newcommand{\pts}[1]{\hfill\textit{[#1]}}

\begin{document}
\begin{center}
  {\large\bfseries BANARAS HINDU UNIVERSITY}\\[2pt]
  {\normalsize B.Sc.\ (Hons.) Semester V Examination, 2022--23}\\[6pt]
  \rule{0.8\linewidth}{0.5pt}\\[6pt]
  {\large\bfseries Paper: CS-106 --- Operating System Concepts}\\[4pt]
  {\normalsize Time: Three Hours \hfill Maximum Marks: 70}
\end{center}

\rule{\linewidth}{0.4pt}

\smallskip
\noindent\textit{Instructions: Answer any \textbf{five} questions. All questions carry equal marks.}

\bigskip

\noindent\textbf{Question 1.}\pts{$4+3+3+4=14$}
\begin{parts}
  \item Explain the different states of a process with the help of a state diagram.
  \item Name any two differences between logical and physical addresses.
  \item What is meant by pre-paging? Is it better than demand paging? Justify.
  \item What will be the output of the following code? Draw the execution tree explaining all process creations.
\begin{lstlisting}
#include <stdio.h>
#include <unistd.h>
int main() {
    if (fork() || fork())
        fork();
    printf("1 ");
    return 0;
}
\end{lstlisting}
\end{parts}

\bigskip

\noindent\textbf{Question 2.}\pts{$2+3+3+3+3=14$}
\begin{parts}
  \item List out the benefits and challenges of thread handling.
  \item What is thrashing? How can this problem be resolved?
  \item Explain the principle of locality.
  \item Explain the working set model.
  \item Explain Compaction with a suitable example.
\end{parts}

\bigskip

\noindent\textbf{Question 3.}\pts{$7+7=14$}
\begin{parts}
  \item Why are Translation Look-aside Buffers (TLBs) important? Explain the details stored in a TLB table entry.
  \item Explain the concept of demand paging in detail with neat diagrams.
\end{parts}

\bigskip

\noindent\textbf{Question 4.}\pts{$4+10=14$}
\begin{parts}
  \item What are the different principles which must be considered while selecting a CPU scheduling algorithm? Support your answer with an example.
  \item Consider the following set of processes with CPU burst times (in milliseconds):

  \begin{center}
  \begin{tabular}{|c|c|c|c|}
  \hline
  \textbf{Process} & \textbf{Arrival Time} & \textbf{Burst Time} & \textbf{Priority} \\
  \hline
  P1 & 0  & 3 & 3 \\
  P2 & 1  & 5 & 1 \\
  P3 & 3  & 2 & 3 \\
  P4 & 9  & 5 & 2 \\
  P5 & 12 & 5 & 1 \\
  \hline
  \end{tabular}
  \end{center}

  Draw a Gantt chart and find the average waiting time and average turnaround time for:
  \begin{enumerate}[label=(\roman*)]
    \item FCFS
    \item SJF (Preemptive)
    \item SJF (Non-Preemptive)
    \item Priority Scheduling (Preemptive; smaller number = higher priority)
    \item Round Robin (Time Quantum = 1 ms)
  \end{enumerate}
\end{parts}

\bigskip

\noindent\textbf{Question 5.}\pts{$6+8=14$}
\begin{parts}
  \item Compare the following memory management techniques on their strengths and weaknesses: (i) Fixed Partition, (ii) Dynamic Partition, (iii) Simple Paging.
  \item Given the following page reference string:\\
  1, 2, 3, 2, 5, 6, 3, 4, 6, 3, 7, 3, 1, 6, 3, 4, 5, 3, 2, 4, 3, 4, 5, 1\\
  Number of page frames = 4. Show the page trace and calculate the number of page faults for the following page replacement policies: (i)~LRU, (ii)~Optimal, (iii)~FIFO. Also explain Belady's anomaly.
\end{parts}

\bigskip

\noindent\textbf{Question 6.}\pts{$8+6=14$}
\begin{parts}
  \item What is busy waiting with respect to process synchronization? Explain how a semaphore reduces the severity of this problem. Also define and distinguish between:
    \begin{enumerate}[label=(\roman*)]
      \item General (counting) semaphores
      \item Binary semaphores
      \item Strong semaphores
      \item Weak semaphores
    \end{enumerate}
  \item The following program consists of 3 concurrent processes (P0, P1, P2) and 3 binary semaphores initialized as S0=1, S1=0, S2=0:
\begin{lstlisting}[numbers=none]
/* P0 */          /* P1 */          /* P2 */
while(true){      while(true){      while(true){
  wait(S0);         wait(S1);         wait(S2);
  print("0");       print("1");       print("2");
  signal(S1);       signal(S0);       signal(S0);
  signal(S2);
}                 }                 }
\end{lstlisting}
  How many times will process P0 print '0'? Give step-by-step explanation.
\end{parts}

\bigskip

\noindent\textbf{Question 7.} Write the Banker's Algorithm. Consider the following snapshot of a system at time $T_0$ with 5 processes (P0--P4) and 4 resource types (A, B, C, D):\pts{$3+4+3+4=14$}

\begin{center}
\begin{tabular}{|c|cccc|cccc|cccc|}
\hline
& \multicolumn{4}{c|}{\textbf{Allocation}} & \multicolumn{4}{c|}{\textbf{Maximum}} & \multicolumn{4}{c|}{\textbf{Available}} \\
\textbf{Process} & A & B & C & D & A & B & C & D & A & B & C & D \\
\hline
P0 & 0 & 0 & 1 & 2 & 0 & 0 & 1 & 2 & 1 & 5 & 2 & 0 \\
P1 & 1 & 0 & 0 & 0 & 1 & 7 & 5 & 0 &   &   &   &   \\
P2 & 1 & 3 & 5 & 4 & 2 & 3 & 5 & 6 &   &   &   &   \\
P3 & 0 & 6 & 3 & 2 & 0 & 6 & 5 & 2 &   &   &   &   \\
P4 & 0 & 0 & 1 & 4 & 0 & 6 & 5 & 6 &   &   &   &   \\
\hline
\end{tabular}
\end{center}

Answer the following:
\begin{parts}
  \item What is the content of the NEED matrix?
  \item Is the system in a safe state? If yes, justify with step-wise calculations.
  \item Which processes may cause deadlock if the system is not safe?
  \item If a request from process P1 arrives for (0, 4, 3, 1) --- can the request be granted immediately? Justify.
\end{parts}

\vfill
\rule{\linewidth}{0.4pt}
\begin{center}\small
\href{https://bhu-exam.vercel.app/}{Click here to visit our website}\\[4pt]\href{https://me-aryan.vercel.app/}{Click here to visit our contributor}\\[4pt]\href{https://quashan-me.vercel.app/}{Click here to visit our contributor}
\end{center}
\end{document}